%%%%%%%%%%%%%%%%%%%%%%%%%%%%%%%%%%%%%%%%%%%%%%%%%%%%%%%%%%%%
% 7.5 JOINT DISTRIBUTIONS
% The Composition of Advanced Mathematics
%%%%%%%%%%%%%%%%%%%%%%%%%%%%%%%%%%%%%%%%%%%%%%%%%%%%%%%%%%%%

\section{Joint Distributions}
\label{sec:joint-distributions}

\prerequisite{
Random variables, discrete and continuous distributions, probability
mass functions, probability density functions, cumulative distribution
functions, conditional probability, independence, double integrals, and
basic multivariable calculus.
}

\learningobjectives{
By the end of this section, the reader should be able to:
\begin{itemize}
    \item define a random vector;
    \item work with joint probability mass functions;
    \item work with joint probability density functions;
    \item verify joint distributions;
    \item compute marginal distributions;
    \item compute conditional distributions;
    \item distinguish joint, marginal, and conditional distributions;
    \item determine independence from a joint distribution;
    \item use joint cumulative distribution functions;
    \item interpret support regions in two dimensions;
    \item compute probabilities over regions of the plane;
    \item use iterated and double integrals for joint densities;
    \item recognize conditional density formulas;
    \item understand products of independent marginal distributions;
    \item work with several discrete and continuous random variables;
    \item understand covariance conceptually as a measure of joint
          behavior;
    \item recognize correlation as a standardized dependence measure;
    \item understand that zero covariance does not generally imply
          independence;
    \item work with conditional independence;
    \item recognize exchangeable random variables;
    \item understand multivariate normal distributions at an
          introductory level;
    \item connect joint distributions with statistics, stochastic
          processes, data analysis, and multivariate probability.
\end{itemize}
}

%%%%%%%%%%%%%%%%%%%%%%%%%%%%%%%%%%%%%%%%%%%%%%%%%%%%%%%%%%%%
\subsection{From One Random Variable to Several}
\label{subsec:from-one-rv-to-several}
%%%%%%%%%%%%%%%%%%%%%%%%%%%%%%%%%%%%%%%%%%%%%%%%%%%%%%%%%%%%

Many experiments naturally produce more than one numerical quantity.

For example, rolling two dice produces two numbers.

A weather observation may record both:

\[
\text{temperature}
\]

and

\[
\text{humidity}.
\]

A statistical sample may contain many random measurements.

To model several random quantities simultaneously, we use random
vectors.

%%%%%%%%%%%%%%%%%%%%%%%%%%%%%%%%%%%%%%%%%%%%%%%%%%%%%%%%%%%%
\subsection{Random Vectors}
\label{subsec:random-vectors}
%%%%%%%%%%%%%%%%%%%%%%%%%%%%%%%%%%%%%%%%%%%%%%%%%%%%%%%%%%%%

\begin{definitionbox}{Random Vector}
A random vector is a vector-valued function

\[
\boxed{
\mathbf{X}
=
(X_1,\ldots,X_n):
\Omega\to\mathbb{R}^n.
}
\]
\end{definitionbox}

Each coordinate

\[
X_i
\]

is itself a random variable.

For two variables, we often write

\[
\boxed{
(X,Y).
}
\]

This is called a bivariate random vector.

%%%%%%%%%%%%%%%%%%%%%%%%%%%%%%%%%%%%%%%%%%%%%%%%%%%%%%%%%%%%
\subsection{Joint Distribution}
\label{subsec:joint-distribution-definition}
%%%%%%%%%%%%%%%%%%%%%%%%%%%%%%%%%%%%%%%%%%%%%%%%%%%%%%%%%%%%

The joint distribution describes the simultaneous probability behavior
of several random variables.

For two random variables \(X\) and \(Y\), it answers questions such as

\[
P(X=x,Y=y),
\]

\[
P(X\leq a,Y\leq b),
\]

and

\[
P(X+Y>c).
\]

Thus,

\[
\boxed{
\text{joint distribution}
=
\text{simultaneous law of several random variables}.
}
\]

%%%%%%%%%%%%%%%%%%%%%%%%%%%%%%%%%%%%%%%%%%%%%%%%%%%%%%%%%%%%
\subsection{Discrete Joint Distributions}
\label{subsec:discrete-joint-distributions}
%%%%%%%%%%%%%%%%%%%%%%%%%%%%%%%%%%%%%%%%%%%%%%%%%%%%%%%%%%%%

Suppose \(X\) and \(Y\) are discrete random variables.

\begin{definitionbox}{Joint Probability Mass Function}
The joint probability mass function of \(X\) and \(Y\) is

\[
\boxed{
p_{X,Y}(x,y)
=
P(X=x,Y=y).
}
\]
\end{definitionbox}

The pair

\[
(x,y)
\]

must belong to the joint support.

%%%%%%%%%%%%%%%%%%%%%%%%%%%%%%%%%%%%%%%%%%%%%%%%%%%%%%%%%%%%
\subsection{Properties of a Joint PMF}
\label{subsec:joint-pmf-properties}
%%%%%%%%%%%%%%%%%%%%%%%%%%%%%%%%%%%%%%%%%%%%%%%%%%%%%%%%%%%%

A joint PMF satisfies

\[
\boxed{
p_{X,Y}(x,y)\geq0
}
\]

for all \(x,y\), and

\[
\boxed{
\sum_x\sum_y
p_{X,Y}(x,y)
=
1.
}
\]

For an event

\[
A\subseteq\mathbb{R}^2,
\]

we have

\[
\boxed{
P((X,Y)\in A)
=
\sum_{(x,y)\in A}
p_{X,Y}(x,y).
}
\]

%%%%%%%%%%%%%%%%%%%%%%%%%%%%%%%%%%%%%%%%%%%%%%%%%%%%%%%%%%%%
\subsection{Joint Probability Tables}
\label{subsec:joint-probability-tables}
%%%%%%%%%%%%%%%%%%%%%%%%%%%%%%%%%%%%%%%%%%%%%%%%%%%%%%%%%%%%

For finite discrete variables, a joint PMF can be represented by a
table.

For example,

\[
\begin{array}{c|ccc|c}
 & Y=0 & Y=1 & Y=2 & p_X(x)\\
\hline
X=0 & 0.10 & 0.20 & 0.10 & 0.40\\
X=1 & 0.15 & 0.25 & 0.20 & 0.60\\
\hline
p_Y(y) & 0.25 & 0.45 & 0.30 & 1
\end{array}
\]

The interior entries form the joint PMF.

The row and column totals give marginal distributions.

%%%%%%%%%%%%%%%%%%%%%%%%%%%%%%%%%%%%%%%%%%%%%%%%%%%%%%%%%%%%
\subsection{Marginal Distributions}
\label{subsec:marginal-distributions}
%%%%%%%%%%%%%%%%%%%%%%%%%%%%%%%%%%%%%%%%%%%%%%%%%%%%%%%%%%%%

The distribution of one coordinate by itself is called a marginal
distribution.

For discrete variables,

\[
\boxed{
p_X(x)
=
\sum_y
p_{X,Y}(x,y).
}
\]

Similarly,

\[
\boxed{
p_Y(y)
=
\sum_x
p_{X,Y}(x,y).
}
\]

The word ``marginal'' comes from the totals often written in the margins
of a probability table.

%%%%%%%%%%%%%%%%%%%%%%%%%%%%%%%%%%%%%%%%%%%%%%%%%%%%%%%%%%%%
\subsection{Worked Example: Finding Marginals}
\label{subsec:worked-discrete-marginals}
%%%%%%%%%%%%%%%%%%%%%%%%%%%%%%%%%%%%%%%%%%%%%%%%%%%%%%%%%%%%

\begin{workedexamplebox}{Marginals from a Joint Table}

Suppose

\[
\begin{array}{c|cc}
 & Y=0 & Y=1\\
\hline
X=0 & 0.20 & 0.30\\
X=1 & 0.10 & 0.40
\end{array}
\]

Then

\[
P(X=0)
=
0.20+0.30
=
0.50,
\]

and

\[
P(X=1)
=
0.10+0.40
=
0.50.
\]

Also,

\[
P(Y=0)
=
0.20+0.10
=
0.30,
\]

and

\[
P(Y=1)
=
0.30+0.40
=
0.70.
\]

\begin{answerbox}
\[
\boxed{
p_X(0)=0.50,\quad
p_X(1)=0.50,
}
\]

\[
\boxed{
p_Y(0)=0.30,\quad
p_Y(1)=0.70.
}
\]
\end{answerbox}

\end{workedexamplebox}

%%%%%%%%%%%%%%%%%%%%%%%%%%%%%%%%%%%%%%%%%%%%%%%%%%%%%%%%%%%%
\subsection{Conditional PMFs}
\label{subsec:conditional-pmfs}
%%%%%%%%%%%%%%%%%%%%%%%%%%%%%%%%%%%%%%%%%%%%%%%%%%%%%%%%%%%%

For discrete variables, the conditional PMF of \(X\) given \(Y=y\) is

\[
\boxed{
p_{X\mid Y}(x\mid y)
=
P(X=x\mid Y=y).
}
\]

If

\[
P(Y=y)>0,
\]

then

\[
\boxed{
p_{X\mid Y}(x\mid y)
=
\frac{
p_{X,Y}(x,y)
}{
p_Y(y)
}.
}
\]

%%%%%%%%%%%%%%%%%%%%%%%%%%%%%%%%%%%%%%%%%%%%%%%%%%%%%%%%%%%%
\subsection{Worked Example: Conditional PMF}
\label{subsec:worked-conditional-pmf}
%%%%%%%%%%%%%%%%%%%%%%%%%%%%%%%%%%%%%%%%%%%%%%%%%%%%%%%%%%%%

\begin{workedexamplebox}{Conditioning on \(Y=1\)}

Use the joint PMF

\[
\begin{array}{c|cc}
 & Y=0 & Y=1\\
\hline
X=0 & 0.20 & 0.30\\
X=1 & 0.10 & 0.40
\end{array}
\]

We found

\[
P(Y=1)=0.70.
\]

Therefore,

\[
P(X=0\mid Y=1)
=
\frac{0.30}{0.70}
=
\frac37,
\]

and

\[
P(X=1\mid Y=1)
=
\frac{0.40}{0.70}
=
\frac47.
\]

\begin{answerbox}
\[
\boxed{
p_{X\mid Y}(0\mid1)=\frac37,
\qquad
p_{X\mid Y}(1\mid1)=\frac47.
}
\]
\end{answerbox}

\end{workedexamplebox}

%%%%%%%%%%%%%%%%%%%%%%%%%%%%%%%%%%%%%%%%%%%%%%%%%%%%%%%%%%%%
\subsection{Discrete Independence}
\label{subsec:discrete-joint-independence}
%%%%%%%%%%%%%%%%%%%%%%%%%%%%%%%%%%%%%%%%%%%%%%%%%%%%%%%%%%%%

Two discrete random variables are independent if

\[
\boxed{
p_{X,Y}(x,y)
=
p_X(x)p_Y(y)
}
\]

for every pair \((x,y)\).

Thus the joint PMF factors into the product of its marginals.

%%%%%%%%%%%%%%%%%%%%%%%%%%%%%%%%%%%%%%%%%%%%%%%%%%%%%%%%%%%%
\subsection{Worked Example: Testing Independence}
\label{subsec:worked-discrete-independence}
%%%%%%%%%%%%%%%%%%%%%%%%%%%%%%%%%%%%%%%%%%%%%%%%%%%%%%%%%%%%

\begin{workedexamplebox}{Checking Factorization}

Suppose

\[
P(X=0)=P(X=1)=\frac12
\]

and

\[
P(Y=0)=\frac13,
\qquad
P(Y=1)=\frac23.
\]

If \(X\) and \(Y\) are independent, then

\[
P(X=1,Y=1)
=
\frac12\cdot\frac23
=
\frac13.
\]

If the actual joint probability is instead

\[
\frac14,
\]

then factorization fails.

\begin{answerbox}
\[
\boxed{
X\text{ and }Y\text{ are not independent}.
}
\]
\end{answerbox}

\end{workedexamplebox}

%%%%%%%%%%%%%%%%%%%%%%%%%%%%%%%%%%%%%%%%%%%%%%%%%%%%%%%%%%%%
\subsection{Joint Support}
\label{subsec:joint-support}
%%%%%%%%%%%%%%%%%%%%%%%%%%%%%%%%%%%%%%%%%%%%%%%%%%%%%%%%%%%%

The joint support is the set of pairs where the joint distribution is
concentrated.

For discrete variables,

\[
\boxed{
\operatorname{supp}(X,Y)
=
\{
(x,y):
p_{X,Y}(x,y)>0
\}.
}
\]

For continuous variables, the support is typically a region in the
plane.

The geometry of this region is often essential.

%%%%%%%%%%%%%%%%%%%%%%%%%%%%%%%%%%%%%%%%%%%%%%%%%%%%%%%%%%%%
\subsection{Continuous Joint Distributions}
\label{subsec:continuous-joint-distributions}
%%%%%%%%%%%%%%%%%%%%%%%%%%%%%%%%%%%%%%%%%%%%%%%%%%%%%%%%%%%%

For continuous random variables \(X\) and \(Y\), probabilities are
described by a joint probability density function.

\begin{definitionbox}{Joint Probability Density Function}
A joint density \(f_{X,Y}\) satisfies

\[
\boxed{
f_{X,Y}(x,y)\geq0
}
\]

and

\[
\boxed{
\iint_{\mathbb{R}^2}
f_{X,Y}(x,y)\,dx\,dy
=
1.
}
\]
\end{definitionbox}

For a region \(A\subseteq\mathbb{R}^2\),

\[
\boxed{
P((X,Y)\in A)
=
\iint_A
f_{X,Y}(x,y)\,dx\,dy.
}
\]

%%%%%%%%%%%%%%%%%%%%%%%%%%%%%%%%%%%%%%%%%%%%%%%%%%%%%%%%%%%%
\subsection{Probability over Rectangles}
\label{subsec:joint-probability-rectangles}
%%%%%%%%%%%%%%%%%%%%%%%%%%%%%%%%%%%%%%%%%%%%%%%%%%%%%%%%%%%%

For a rectangle

\[
a<X<b,
\qquad
c<Y<d,
\]

we have

\[
\boxed{
P(a<X<b,\;c<Y<d)
=
\int_a^b
\int_c^d
f_{X,Y}(x,y)\,dy\,dx.
}
\]

The order of integration may often be reversed when the relevant
conditions for Fubini's theorem are satisfied.

%%%%%%%%%%%%%%%%%%%%%%%%%%%%%%%%%%%%%%%%%%%%%%%%%%%%%%%%%%%%
\subsection{Worked Example: Constant Joint Density}
\label{subsec:worked-constant-joint-density}
%%%%%%%%%%%%%%%%%%%%%%%%%%%%%%%%%%%%%%%%%%%%%%%%%%%%%%%%%%%%

\begin{workedexamplebox}{Uniform Density on a Rectangle}

Suppose

\[
f_{X,Y}(x,y)
=
c
\]

on

\[
0\leq x\leq2,
\qquad
0\leq y\leq3,
\]

and zero elsewhere.

Normalization requires

\[
\int_0^2\int_0^3 c\,dy\,dx
=
1.
\]

Thus,

\[
6c=1.
\]

Therefore,

\begin{answerbox}
\[
\boxed{
c=\frac16.
}
\]
\end{answerbox}

\end{workedexamplebox}

%%%%%%%%%%%%%%%%%%%%%%%%%%%%%%%%%%%%%%%%%%%%%%%%%%%%%%%%%%%%
\subsection{Continuous Marginal Densities}
\label{subsec:continuous-marginal-densities}
%%%%%%%%%%%%%%%%%%%%%%%%%%%%%%%%%%%%%%%%%%%%%%%%%%%%%%%%%%%%

The marginal density of \(X\) is obtained by integrating out \(Y\):

\[
\boxed{
f_X(x)
=
\int_{-\infty}^{\infty}
f_{X,Y}(x,y)\,dy.
}
\]

Similarly,

\[
\boxed{
f_Y(y)
=
\int_{-\infty}^{\infty}
f_{X,Y}(x,y)\,dx.
}
\]

The phrase ``integrating out'' means removing one variable by summing
its density contribution across all possible values.

%%%%%%%%%%%%%%%%%%%%%%%%%%%%%%%%%%%%%%%%%%%%%%%%%%%%%%%%%%%%
\subsection{Worked Example: Marginals from a Joint Density}
\label{subsec:worked-continuous-marginals}
%%%%%%%%%%%%%%%%%%%%%%%%%%%%%%%%%%%%%%%%%%%%%%%%%%%%%%%%%%%%

\begin{workedexamplebox}{Uniform Rectangle Marginals}

Continue with

\[
f_{X,Y}(x,y)
=
\frac16
\]

for

\[
0\leq x\leq2,
\qquad
0\leq y\leq3.
\]

Then for \(0\leq x\leq2\),

\[
f_X(x)
=
\int_0^3
\frac16\,dy
=
\frac12.
\]

For \(0\leq y\leq3\),

\[
f_Y(y)
=
\int_0^2
\frac16\,dx
=
\frac13.
\]

\begin{answerbox}
\[
\boxed{
X\sim\operatorname{Uniform}(0,2),
}
\]

\[
\boxed{
Y\sim\operatorname{Uniform}(0,3).
}
\]
\end{answerbox}

\end{workedexamplebox}

%%%%%%%%%%%%%%%%%%%%%%%%%%%%%%%%%%%%%%%%%%%%%%%%%%%%%%%%%%%%
\subsection{Continuous Conditional Densities}
\label{subsec:continuous-conditional-densities}
%%%%%%%%%%%%%%%%%%%%%%%%%%%%%%%%%%%%%%%%%%%%%%%%%%%%%%%%%%%%

If

\[
f_Y(y)>0,
\]

the conditional density of \(X\) given \(Y=y\) is

\[
\boxed{
f_{X\mid Y}(x\mid y)
=
\frac{
f_{X,Y}(x,y)
}{
f_Y(y)
}.
}
\]

Similarly,

\[
\boxed{
f_{Y\mid X}(y\mid x)
=
\frac{
f_{X,Y}(x,y)
}{
f_X(x)
}.
}
\]

%%%%%%%%%%%%%%%%%%%%%%%%%%%%%%%%%%%%%%%%%%%%%%%%%%%%%%%%%%%%
\subsection{Conditional Density Interpretation}
\label{subsec:conditional-density-interpretation}
%%%%%%%%%%%%%%%%%%%%%%%%%%%%%%%%%%%%%%%%%%%%%%%%%%%%%%%%%%%%

A conditional density describes the probability behavior of one
variable when the other is fixed.

For a fixed value \(y\),

\[
f_{X\mid Y}(x\mid y)
\]

must satisfy

\[
\boxed{
\int_{-\infty}^{\infty}
f_{X\mid Y}(x\mid y)\,dx
=
1.
}
\]

Thus each fixed \(y\) produces a full probability density in \(x\).

%%%%%%%%%%%%%%%%%%%%%%%%%%%%%%%%%%%%%%%%%%%%%%%%%%%%%%%%%%%%
\subsection{Worked Example: Conditional Density}
\label{subsec:worked-conditional-density}
%%%%%%%%%%%%%%%%%%%%%%%%%%%%%%%%%%%%%%%%%%%%%%%%%%%%%%%%%%%%

\begin{workedexamplebox}{Conditioning on a Uniform Rectangle}

Again suppose

\[
f_{X,Y}(x,y)
=
\frac16
\]

on

\[
0\leq x\leq2,
\qquad
0\leq y\leq3.
\]

We have

\[
f_Y(y)=\frac13
\]

for

\[
0\leq y\leq3.
\]

Therefore,

\[
f_{X\mid Y}(x\mid y)
=
\frac{1/6}{1/3}
=
\frac12,
\]

for

\[
0\leq x\leq2.
\]

\begin{answerbox}
\[
\boxed{
X\mid(Y=y)
\sim
\operatorname{Uniform}(0,2)
}
\]

for every admissible \(y\).
\end{answerbox}

\end{workedexamplebox}

%%%%%%%%%%%%%%%%%%%%%%%%%%%%%%%%%%%%%%%%%%%%%%%%%%%%%%%%%%%%
\subsection{Continuous Independence}
\label{subsec:continuous-independence}
%%%%%%%%%%%%%%%%%%%%%%%%%%%%%%%%%%%%%%%%%%%%%%%%%%%%%%%%%%%%

Continuous random variables \(X\) and \(Y\) are independent if

\[
\boxed{
f_{X,Y}(x,y)
=
f_X(x)f_Y(y)
}
\]

for all relevant \(x,y\).

Thus independence again corresponds to factorization of the joint law.

%%%%%%%%%%%%%%%%%%%%%%%%%%%%%%%%%%%%%%%%%%%%%%%%%%%%%%%%%%%%
\subsection{Rectangular Support and Independence}
\label{subsec:rectangular-support-independence}
%%%%%%%%%%%%%%%%%%%%%%%%%%%%%%%%%%%%%%%%%%%%%%%%%%%%%%%%%%%%

A factorized joint density typically has support of product form:

\[
\boxed{
S_X\times S_Y.
}
\]

However, rectangular support alone does not guarantee independence.

One must also verify factorization of the density.

Conversely, a nonrectangular support often immediately proves
dependence.

%%%%%%%%%%%%%%%%%%%%%%%%%%%%%%%%%%%%%%%%%%%%%%%%%%%%%%%%%%%%
\subsection{Worked Example: Nonrectangular Support}
\label{subsec:worked-nonrectangular-support}
%%%%%%%%%%%%%%%%%%%%%%%%%%%%%%%%%%%%%%%%%%%%%%%%%%%%%%%%%%%%

\begin{workedexamplebox}{Triangular Support}

Suppose

\[
f_{X,Y}(x,y)
=
2
\]

on the region

\[
0<y<x<1,
\]

and zero elsewhere.

The support is triangular.

If \(X\) and \(Y\) were independent, the support would need to factor
into independent allowable ranges for \(X\) and \(Y\).

But the restriction

\[
y<x
\]

links the variables directly.

\begin{answerbox}
\[
\boxed{
X\text{ and }Y\text{ are not independent}.
}
\]
\end{answerbox}

\end{workedexamplebox}

%%%%%%%%%%%%%%%%%%%%%%%%%%%%%%%%%%%%%%%%%%%%%%%%%%%%%%%%%%%%
\subsection{Probability over General Regions}
\label{subsec:probability-general-regions}
%%%%%%%%%%%%%%%%%%%%%%%%%%%%%%%%%%%%%%%%%%%%%%%%%%%%%%%%%%%%

Suppose \(A\) is described by

\[
a\leq x\leq b
\]

and

\[
g_1(x)\leq y\leq g_2(x).
\]

Then

\[
\boxed{
P((X,Y)\in A)
=
\int_a^b
\int_{g_1(x)}^{g_2(x)}
f_{X,Y}(x,y)\,dy\,dx.
}
\]

Choosing correct integration limits is often the main challenge.

%%%%%%%%%%%%%%%%%%%%%%%%%%%%%%%%%%%%%%%%%%%%%%%%%%%%%%%%%%%%
\subsection{Worked Example: Probability over a Triangle}
\label{subsec:worked-triangular-region-probability}
%%%%%%%%%%%%%%%%%%%%%%%%%%%%%%%%%%%%%%%%%%%%%%%%%%%%%%%%%%%%

\begin{workedexamplebox}{Integrating over a Subregion}

Suppose

\[
f_{X,Y}(x,y)=2
\]

for

\[
0<y<x<1.
\]

Find

\[
P(X<1/2).
\]

The appropriate region is

\[
0<x<\frac12,
\qquad
0<y<x.
\]

Thus,

\[
P(X<1/2)
=
\int_0^{1/2}
\int_0^x
2\,dy\,dx.
\]

Integrating,

\[
=
\int_0^{1/2}
2x\,dx
=
\left[
x^2
\right]_0^{1/2}.
\]

\begin{answerbox}
\[
\boxed{
P(X<1/2)
=
\frac14.
}
\]
\end{answerbox}

\end{workedexamplebox}

%%%%%%%%%%%%%%%%%%%%%%%%%%%%%%%%%%%%%%%%%%%%%%%%%%%%%%%%%%%%
\subsection{Joint Cumulative Distribution Function}
\label{subsec:joint-cdf}
%%%%%%%%%%%%%%%%%%%%%%%%%%%%%%%%%%%%%%%%%%%%%%%%%%%%%%%%%%%%

\begin{definitionbox}{Joint CDF}
The joint cumulative distribution function of \(X\) and \(Y\) is

\[
\boxed{
F_{X,Y}(x,y)
=
P(X\leq x,\;Y\leq y).
}
\]
\end{definitionbox}

For continuous variables,

\[
\boxed{
F_{X,Y}(x,y)
=
\int_{-\infty}^{x}
\int_{-\infty}^{y}
f_{X,Y}(u,v)\,dv\,du.
}
\]

%%%%%%%%%%%%%%%%%%%%%%%%%%%%%%%%%%%%%%%%%%%%%%%%%%%%%%%%%%%%
\subsection{Recovering Marginal CDFs}
\label{subsec:marginals-from-joint-cdf}
%%%%%%%%%%%%%%%%%%%%%%%%%%%%%%%%%%%%%%%%%%%%%%%%%%%%%%%%%%%%

The marginal CDF of \(X\) is obtained by allowing \(Y\) to range over all
values:

\[
\boxed{
F_X(x)
=
\lim_{y\to\infty}
F_{X,Y}(x,y).
}
\]

Similarly,

\[
\boxed{
F_Y(y)
=
\lim_{x\to\infty}
F_{X,Y}(x,y).
}
\]

Thus the joint CDF contains the marginal distributions.

%%%%%%%%%%%%%%%%%%%%%%%%%%%%%%%%%%%%%%%%%%%%%%%%%%%%%%%%%%%%
\subsection{Rectangle Probabilities from a Joint CDF}
\label{subsec:rectangle-probabilities-joint-cdf}
%%%%%%%%%%%%%%%%%%%%%%%%%%%%%%%%%%%%%%%%%%%%%%%%%%%%%%%%%%%%

For

\[
a<b
\]

and

\[
c<d,
\]

we have

\[
\boxed{
\begin{aligned}
P(a<X\leq b,\;c<Y\leq d)
={}&
F_{X,Y}(b,d)
-F_{X,Y}(a,d)
\\
&-
F_{X,Y}(b,c)
+
F_{X,Y}(a,c).
\end{aligned}
}
\]

This is a two-dimensional inclusion--exclusion formula.

%%%%%%%%%%%%%%%%%%%%%%%%%%%%%%%%%%%%%%%%%%%%%%%%%%%%%%%%%%%%
\subsection{Joint Density from a Joint CDF}
\label{subsec:joint-density-from-cdf}
%%%%%%%%%%%%%%%%%%%%%%%%%%%%%%%%%%%%%%%%%%%%%%%%%%%%%%%%%%%%

If the joint CDF is sufficiently differentiable, then

\[
\boxed{
f_{X,Y}(x,y)
=
\frac{\partial^2}{
\partial x\,\partial y
}
F_{X,Y}(x,y).
}
\]

Thus the joint density is obtained by differentiating once with respect
to each variable.

%%%%%%%%%%%%%%%%%%%%%%%%%%%%%%%%%%%%%%%%%%%%%%%%%%%%%%%%%%%%
\subsection{Independence via CDFs}
\label{subsec:independence-via-cdfs}
%%%%%%%%%%%%%%%%%%%%%%%%%%%%%%%%%%%%%%%%%%%%%%%%%%%%%%%%%%%%

Random variables \(X\) and \(Y\) are independent if and only if

\[
\boxed{
F_{X,Y}(x,y)
=
F_X(x)F_Y(y)
}
\]

for all \(x,y\).

This criterion works beyond purely discrete or absolutely continuous
cases.

%%%%%%%%%%%%%%%%%%%%%%%%%%%%%%%%%%%%%%%%%%%%%%%%%%%%%%%%%%%%
\subsection{Functions of Two Random Variables}
\label{subsec:functions-two-random-variables}
%%%%%%%%%%%%%%%%%%%%%%%%%%%%%%%%%%%%%%%%%%%%%%%%%%%%%%%%%%%%

A function of several random variables produces a new random variable.

For example,

\[
\boxed{
Z=g(X,Y).
}
\]

Common examples include

\[
Z=X+Y,
\]

\[
Z=XY,
\]

\[
Z=\max(X,Y),
\]

and

\[
Z=\min(X,Y).
\]

Their distributions will be developed systematically in
Section~\ref{sec:functions-random-variables}.

%%%%%%%%%%%%%%%%%%%%%%%%%%%%%%%%%%%%%%%%%%%%%%%%%%%%%%%%%%%%
\subsection{Distribution of a Sum}
\label{subsec:distribution-sum-preview}
%%%%%%%%%%%%%%%%%%%%%%%%%%%%%%%%%%%%%%%%%%%%%%%%%%%%%%%%%%%%

For discrete independent variables,

\[
\boxed{
P(X+Y=z)
=
\sum_x
P(X=x)P(Y=z-x).
}
\]

This is a discrete convolution.

For independent continuous variables,

\[
\boxed{
f_{X+Y}(z)
=
\int_{-\infty}^{\infty}
f_X(x)f_Y(z-x)\,dx.
}
\]

This is continuous convolution.

%%%%%%%%%%%%%%%%%%%%%%%%%%%%%%%%%%%%%%%%%%%%%%%%%%%%%%%%%%%%
\subsection{Worked Example: Sum of Two Dice}
\label{subsec:worked-sum-two-dice}
%%%%%%%%%%%%%%%%%%%%%%%%%%%%%%%%%%%%%%%%%%%%%%%%%%%%%%%%%%%%

\begin{workedexamplebox}{Two Fair Dice}

Let \(X\) and \(Y\) be independent fair die results.

Find

\[
P(X+Y=7).
\]

The favorable ordered pairs are

\[
(1,6),
(2,5),
(3,4),
(4,3),
(5,2),
(6,1).
\]

There are

\[
36
\]

equally likely ordered pairs in total.

Therefore,

\begin{answerbox}
\[
\boxed{
P(X+Y=7)
=
\frac6{36}
=
\frac16.
}
\]
\end{answerbox}

\end{workedexamplebox}

%%%%%%%%%%%%%%%%%%%%%%%%%%%%%%%%%%%%%%%%%%%%%%%%%%%%%%%%%%%%
\subsection{Maximum and Minimum}
\label{subsec:max-min-joint-preview}
%%%%%%%%%%%%%%%%%%%%%%%%%%%%%%%%%%%%%%%%%%%%%%%%%%%%%%%%%%%%

Let

\[
M=\max(X,Y).
\]

Then

\[
\{M\leq m\}
=
\{X\leq m,\;Y\leq m\}.
\]

Therefore,

\[
\boxed{
F_M(m)
=
F_{X,Y}(m,m).
}
\]

If \(X\) and \(Y\) are independent,

\[
\boxed{
F_M(m)
=
F_X(m)F_Y(m).
}
\]

Similarly, for

\[
L=\min(X,Y),
\]

\[
\boxed{
P(L>l)
=
P(X>l,Y>l).
}
\]

%%%%%%%%%%%%%%%%%%%%%%%%%%%%%%%%%%%%%%%%%%%%%%%%%%%%%%%%%%%%
\subsection{Worked Example: Maximum of Independent Uniforms}
\label{subsec:worked-max-uniforms}
%%%%%%%%%%%%%%%%%%%%%%%%%%%%%%%%%%%%%%%%%%%%%%%%%%%%%%%%%%%%

\begin{workedexamplebox}{Maximum of Two Uniform Variables}

Let \(X\) and \(Y\) be independent and uniformly distributed on

\[
[0,1].
\]

Let

\[
M=\max(X,Y).
\]

For

\[
0\leq m\leq1,
\]

\[
F_M(m)
=
P(X\leq m,Y\leq m).
\]

By independence,

\[
F_M(m)
=
m^2.
\]

Differentiating,

\[
f_M(m)=2m.
\]

\begin{answerbox}
\[
\boxed{
f_M(m)=2m,
\qquad
0\leq m\leq1.
}
\]
\end{answerbox}

\end{workedexamplebox}

%%%%%%%%%%%%%%%%%%%%%%%%%%%%%%%%%%%%%%%%%%%%%%%%%%%%%%%%%%%%
\subsection{Covariance Preview}
\label{subsec:covariance-preview-joint}
%%%%%%%%%%%%%%%%%%%%%%%%%%%%%%%%%%%%%%%%%%%%%%%%%%%%%%%%%%%%

Joint distributions allow us to study whether two variables tend to move
together.

The covariance is defined by

\[
\boxed{
\operatorname{Cov}(X,Y)
=
\mathbb{E}
\left[
(X-\mathbb{E}[X])
(Y-\mathbb{E}[Y])
\right].
}
\]

Expectation will be developed fully in the next section.

Covariance measures linear joint variation.

%%%%%%%%%%%%%%%%%%%%%%%%%%%%%%%%%%%%%%%%%%%%%%%%%%%%%%%%%%%%
\subsection{Equivalent Covariance Formula Preview}
\label{subsec:covariance-equivalent-preview}
%%%%%%%%%%%%%%%%%%%%%%%%%%%%%%%%%%%%%%%%%%%%%%%%%%%%%%%%%%%%

When the required expectations exist,

\[
\boxed{
\operatorname{Cov}(X,Y)
=
\mathbb{E}[XY]
-
\mathbb{E}[X]\mathbb{E}[Y].
}
\]

If \(X\) and \(Y\) are independent, then

\[
\mathbb{E}[XY]
=
\mathbb{E}[X]\mathbb{E}[Y],
\]

so

\[
\boxed{
\operatorname{Cov}(X,Y)=0.
}
\]

The converse is not generally true.

%%%%%%%%%%%%%%%%%%%%%%%%%%%%%%%%%%%%%%%%%%%%%%%%%%%%%%%%%%%%
\subsection{Uncorrelated Does Not Mean Independent}
\label{subsec:uncorrelated-not-independent}
%%%%%%%%%%%%%%%%%%%%%%%%%%%%%%%%%%%%%%%%%%%%%%%%%%%%%%%%%%%%

Consider a random variable \(X\) symmetric about \(0\), and define

\[
Y=X^2.
\]

Then \(Y\) is completely determined by \(X\), so the variables are not
independent.

Yet in many symmetric settings,

\[
\mathbb{E}[X]=0
\]

and

\[
\mathbb{E}[X^3]=0,
\]

which can give

\[
\operatorname{Cov}(X,Y)=0.
\]

Thus,

\[
\boxed{
\text{zero covariance}
\not\Longrightarrow
\text{independence}.
}
\]

%%%%%%%%%%%%%%%%%%%%%%%%%%%%%%%%%%%%%%%%%%%%%%%%%%%%%%%%%%%%
\subsection{Correlation Preview}
\label{subsec:correlation-preview}
%%%%%%%%%%%%%%%%%%%%%%%%%%%%%%%%%%%%%%%%%%%%%%%%%%%%%%%%%%%%

The correlation coefficient is

\[
\boxed{
\rho_{X,Y}
=
\frac{
\operatorname{Cov}(X,Y)
}{
\sigma_X\sigma_Y
},
}
\]

when both standard deviations are positive.

The Greek letter

\[
\rho
\]

is called ``rho.''

Correlation is a dimensionless measure of linear association.

%%%%%%%%%%%%%%%%%%%%%%%%%%%%%%%%%%%%%%%%%%%%%%%%%%%%%%%%%%%%
\subsection{Correlation Bounds}
\label{subsec:correlation-bounds}
%%%%%%%%%%%%%%%%%%%%%%%%%%%%%%%%%%%%%%%%%%%%%%%%%%%%%%%%%%%%

When defined,

\[
\boxed{
-1
\leq
\rho_{X,Y}
\leq
1.
}
\]

The extreme values correspond to exact linear relationships under
appropriate nondegeneracy conditions.

These facts will be developed in the section on expectation and moments.

%%%%%%%%%%%%%%%%%%%%%%%%%%%%%%%%%%%%%%%%%%%%%%%%%%%%%%%%%%%%
\subsection{Conditional Independence}
\label{subsec:conditional-independence}
%%%%%%%%%%%%%%%%%%%%%%%%%%%%%%%%%%%%%%%%%%%%%%%%%%%%%%%%%%%%

Random variables \(X\) and \(Y\) may be independent after conditioning
on another variable \(Z\).

One notation is

\[
\boxed{
X\perp Y\mid Z.
}
\]

This means that once \(Z\) is known, additional knowledge of \(X\) gives
no further information about \(Y\), and vice versa.

For discrete variables, a factorization form is

\[
\boxed{
P(X=x,Y=y\mid Z=z)
=
P(X=x\mid Z=z)
P(Y=y\mid Z=z).
}
\]

%%%%%%%%%%%%%%%%%%%%%%%%%%%%%%%%%%%%%%%%%%%%%%%%%%%%%%%%%%%%
\subsection{Conditional Independence Does Not Imply Independence}
\label{subsec:conditional-independence-not-independence}
%%%%%%%%%%%%%%%%%%%%%%%%%%%%%%%%%%%%%%%%%%%%%%%%%%%%%%%%%%%%

Variables may become independent after conditioning even though they are
dependent marginally.

Conversely, variables that are independent marginally may become
dependent after conditioning.

Thus:

\[
\boxed{
\text{marginal independence}
}
\]

and

\[
\boxed{
\text{conditional independence}
}
\]

are distinct concepts.

This distinction is fundamental in statistics and graphical models.

%%%%%%%%%%%%%%%%%%%%%%%%%%%%%%%%%%%%%%%%%%%%%%%%%%%%%%%%%%%%
\subsection{Worked Example: Common Cause Dependence}
\label{subsec:worked-common-cause}
%%%%%%%%%%%%%%%%%%%%%%%%%%%%%%%%%%%%%%%%%%%%%%%%%%%%%%%%%%%%

\begin{workedexamplebox}{Dependence Through a Hidden Variable}

Suppose \(Z\) represents weather conditions.

Let \(X\) represent umbrella use and \(Y\) represent wet roads.

Both depend strongly on \(Z\).

Without conditioning, \(X\) and \(Y\) may be associated.

After conditioning on the weather state \(Z\), the remaining dependence
may disappear in an appropriate model.

\begin{answerbox}
\[
\boxed{
X\perp Y\mid Z
}
\]

can hold even when \(X\) and \(Y\) are not marginally independent.
\end{answerbox}

\end{workedexamplebox}

%%%%%%%%%%%%%%%%%%%%%%%%%%%%%%%%%%%%%%%%%%%%%%%%%%%%%%%%%%%%
\subsection{Exchangeable Random Variables}
\label{subsec:exchangeable-random-variables}
%%%%%%%%%%%%%%%%%%%%%%%%%%%%%%%%%%%%%%%%%%%%%%%%%%%%%%%%%%%%

A collection

\[
X_1,\ldots,X_n
\]

is exchangeable if its joint distribution is unchanged by permuting the
coordinates.

For example,

\[
(X_1,X_2,X_3)
\]

and

\[
(X_3,X_1,X_2)
\]

have the same joint distribution.

Thus the labels do not affect the probabilistic law.

%%%%%%%%%%%%%%%%%%%%%%%%%%%%%%%%%%%%%%%%%%%%%%%%%%%%%%%%%%%%
\subsection{Independent and Identically Distributed Variables}
\label{subsec:iid-random-variables}
%%%%%%%%%%%%%%%%%%%%%%%%%%%%%%%%%%%%%%%%%%%%%%%%%%%%%%%%%%%%

A collection of random variables is independent and identically
distributed when:

\begin{itemize}
    \item the variables are mutually independent;
    \item each has the same distribution.
\end{itemize}

The abbreviation is

\[
\boxed{
\text{i.i.d.}
}
\]

for ``independent and identically distributed.''

Many major probability theorems assume i.i.d. sequences.

%%%%%%%%%%%%%%%%%%%%%%%%%%%%%%%%%%%%%%%%%%%%%%%%%%%%%%%%%%%%
\subsection{i.i.d. Implies Exchangeable}
\label{subsec:iid-implies-exchangeable}
%%%%%%%%%%%%%%%%%%%%%%%%%%%%%%%%%%%%%%%%%%%%%%%%%%%%%%%%%%%%

If

\[
X_1,\ldots,X_n
\]

are i.i.d., then their joint distribution factors:

\[
\boxed{
f(x_1,\ldots,x_n)
=
\prod_{i=1}^{n}f(x_i)
}
\]

in the continuous case.

Permuting coordinates leaves the product unchanged.

Therefore,

\[
\boxed{
\text{i.i.d.}
\Longrightarrow
\text{exchangeable}.
}
\]

The converse is not generally true.

%%%%%%%%%%%%%%%%%%%%%%%%%%%%%%%%%%%%%%%%%%%%%%%%%%%%%%%%%%%%
\subsection{Joint Distributions of More Than Two Variables}
\label{subsec:multivariate-joint-distributions}
%%%%%%%%%%%%%%%%%%%%%%%%%%%%%%%%%%%%%%%%%%%%%%%%%%%%%%%%%%%%

For a discrete random vector

\[
\mathbf{X}
=
(X_1,\ldots,X_n),
\]

the joint PMF is

\[
\boxed{
p_{\mathbf{X}}
(x_1,\ldots,x_n)
=
P(X_1=x_1,\ldots,X_n=x_n).
}
\]

For a continuous random vector,

\[
\boxed{
f_{\mathbf{X}}
(x_1,\ldots,x_n)
}
\]

is a joint density satisfying

\[
\boxed{
\int_{\mathbb{R}^n}
f_{\mathbf{X}}(\mathbf{x})\,d\mathbf{x}
=
1.
}
\]

%%%%%%%%%%%%%%%%%%%%%%%%%%%%%%%%%%%%%%%%%%%%%%%%%%%%%%%%%%%%
\subsection{Marginalization in Higher Dimensions}
\label{subsec:higher-dimensional-marginalization}
%%%%%%%%%%%%%%%%%%%%%%%%%%%%%%%%%%%%%%%%%%%%%%%%%%%%%%%%%%%%

To obtain a marginal distribution, sum or integrate over unwanted
coordinates.

For example,

\[
f_{X_1,X_2}(x_1,x_2)
=
\int_{\mathbb{R}^{n-2}}
f_{\mathbf{X}}
(x_1,\ldots,x_n)
\,dx_3\cdots dx_n.
\]

Thus marginalization systematically reduces dimension.

%%%%%%%%%%%%%%%%%%%%%%%%%%%%%%%%%%%%%%%%%%%%%%%%%%%%%%%%%%%%
\subsection{Multivariate Normal Distribution Preview}
\label{subsec:multivariate-normal-preview}
%%%%%%%%%%%%%%%%%%%%%%%%%%%%%%%%%%%%%%%%%%%%%%%%%%%%%%%%%%%%

The multivariate normal distribution generalizes the ordinary normal
distribution to random vectors.

A multivariate normal vector is determined by:

\[
\boxed{
\text{a mean vector}
}
\]

and

\[
\boxed{
\text{a covariance matrix}.
}
\]

One notation is

\[
\boxed{
\mathbf{X}
\sim
N_n(\boldsymbol{\mu},\Sigma).
}
\]

Here,

\[
\boldsymbol{\mu}
\]

is the mean vector and

\[
\Sigma
\]

is the covariance matrix.

%%%%%%%%%%%%%%%%%%%%%%%%%%%%%%%%%%%%%%%%%%%%%%%%%%%%%%%%%%%%
\subsection{Bivariate Normal Density}
\label{subsec:bivariate-normal-density-preview}
%%%%%%%%%%%%%%%%%%%%%%%%%%%%%%%%%%%%%%%%%%%%%%%%%%%%%%%%%%%%

For standardized variables with correlation \(\rho\), the bivariate
normal density has the form

\[
\boxed{
f(x,y)
=
\frac{
1
}{
2\pi\sqrt{1-\rho^2}
}
\exp
\left[
-
\frac{
x^2-2\rho xy+y^2
}{
2(1-\rho^2)
}
\right].
}
\]

The parameter \(\rho\) controls linear dependence.

When

\[
\rho=0,
\]

the density factors into two independent standard normal densities.

%%%%%%%%%%%%%%%%%%%%%%%%%%%%%%%%%%%%%%%%%%%%%%%%%%%%%%%%%%%%
\subsection{Normality and Zero Correlation}
\label{subsec:normal-zero-correlation}
%%%%%%%%%%%%%%%%%%%%%%%%%%%%%%%%%%%%%%%%%%%%%%%%%%%%%%%%%%%%

For general random variables,

\[
\operatorname{Cov}(X,Y)=0
\]

does not imply independence.

However, for jointly normal random variables,

\[
\boxed{
\operatorname{Cov}(X,Y)=0
\Longrightarrow
X\text{ and }Y\text{ are independent}.
}
\]

This is a special and important property of the multivariate normal
family.

%%%%%%%%%%%%%%%%%%%%%%%%%%%%%%%%%%%%%%%%%%%%%%%%%%%%%%%%%%%%
\subsection{Mixture Distributions Preview}
\label{subsec:mixture-distributions-preview}
%%%%%%%%%%%%%%%%%%%%%%%%%%%%%%%%%%%%%%%%%%%%%%%%%%%%%%%%%%%%

Suppose a latent variable \(Z\) selects among several probability
models.

Then the distribution of \(X\) may be written through the law of total
probability:

\[
\boxed{
P(X\in A)
=
\sum_z
P(X\in A\mid Z=z)P(Z=z)
}
\]

for discrete \(Z\).

Such distributions are called mixtures.

They are central in statistics and machine learning.

%%%%%%%%%%%%%%%%%%%%%%%%%%%%%%%%%%%%%%%%%%%%%%%%%%%%%%%%%%%%
\subsection{Hierarchical Probability Models}
\label{subsec:hierarchical-probability-models}
%%%%%%%%%%%%%%%%%%%%%%%%%%%%%%%%%%%%%%%%%%%%%%%%%%%%%%%%%%%%

A hierarchical model specifies random quantities in stages.

For example,

\[
Z
\sim
\text{some distribution},
\]

then conditionally,

\[
X\mid Z
\sim
\text{another distribution}.
\]

The marginal distribution of \(X\) is obtained by integrating or summing
over \(Z\).

Thus,

\[
\boxed{
\text{conditional model}
+
\text{mixing distribution}
\longrightarrow
\text{marginal model}.
}
\]

%%%%%%%%%%%%%%%%%%%%%%%%%%%%%%%%%%%%%%%%%%%%%%%%%%%%%%%%%%%%
\subsection{Worked Example: Bernoulli Mixture}
\label{subsec:worked-bernoulli-mixture}
%%%%%%%%%%%%%%%%%%%%%%%%%%%%%%%%%%%%%%%%%%%%%%%%%%%%%%%%%%%%

\begin{workedexamplebox}{Random Success Probability}

Suppose

\[
P(Z=1)=0.4,
\qquad
P(Z=2)=0.6.
\]

Conditional on \(Z\),

\[
P(X=1\mid Z=1)=0.2,
\]

and

\[
P(X=1\mid Z=2)=0.8.
\]

By the law of total probability,

\[
P(X=1)
=
(0.2)(0.4)
+
(0.8)(0.6).
\]

Thus,

\begin{answerbox}
\[
\boxed{
P(X=1)=0.56.
}
\]
\end{answerbox}

\end{workedexamplebox}

%%%%%%%%%%%%%%%%%%%%%%%%%%%%%%%%%%%%%%%%%%%%%%%%%%%%%%%%%%%%
\subsection{Joint Distribution and Dependence Structure}
\label{subsec:joint-distribution-dependence}
%%%%%%%%%%%%%%%%%%%%%%%%%%%%%%%%%%%%%%%%%%%%%%%%%%%%%%%%%%%%

Marginal distributions alone do not determine the joint distribution.

Two pairs of random variables may have identical marginals but very
different dependence.

Thus,

\[
\boxed{
\text{marginals}
+
\text{dependence structure}
=
\text{joint distribution}.
}
\]

This distinction becomes especially important in multivariate
statistics.

%%%%%%%%%%%%%%%%%%%%%%%%%%%%%%%%%%%%%%%%%%%%%%%%%%%%%%%%%%%%
\subsection{Worked Example: Same Marginals, Different Dependence}
\label{subsec:worked-same-marginals-different-dependence}
%%%%%%%%%%%%%%%%%%%%%%%%%%%%%%%%%%%%%%%%%%%%%%%%%%%%%%%%%%%%

\begin{workedexamplebox}{Perfect Dependence Versus Independence}

Let \(X\) be Bernoulli with

\[
P(X=0)=P(X=1)=\frac12.
\]

Case 1:

\[
Y=X.
\]

Then \(Y\) has the same marginal distribution as \(X\), but the variables
are perfectly dependent.

Case 2:

Let \(Y\) be an independent Bernoulli\((1/2)\) variable.

Again \(Y\) has the same marginal distribution.

However, the dependence structure is different.

\begin{answerbox}
\[
\boxed{
\text{Marginal distributions do not determine the joint law.}
}
\]
\end{answerbox}

\end{workedexamplebox}

%%%%%%%%%%%%%%%%%%%%%%%%%%%%%%%%%%%%%%%%%%%%%%%%%%%%%%%%%%%%
\subsection{Copula Preview}
\label{subsec:copula-preview}
%%%%%%%%%%%%%%%%%%%%%%%%%%%%%%%%%%%%%%%%%%%%%%%%%%%%%%%%%%%%

A copula is a function that combines marginal distributions into a joint
distribution.

Informally,

\[
\boxed{
\text{copula}
=
\text{mathematical representation of dependence}.
}
\]

Sklar's theorem states that under broad conditions a multivariate
distribution can be written in terms of its marginals and a copula.

This topic is especially important in multivariate statistics and
finance.

%%%%%%%%%%%%%%%%%%%%%%%%%%%%%%%%%%%%%%%%%%%%%%%%%%%%%%%%%%%%
\subsection{Joint Probability Workflow}
\label{subsec:joint-probability-workflow}
%%%%%%%%%%%%%%%%%%%%%%%%%%%%%%%%%%%%%%%%%%%%%%%%%%%%%%%%%%%%

A useful workflow is:

\begin{enumerate}
    \item determine whether the variables are discrete or continuous;
    \item identify the joint support;
    \item verify the joint PMF or density is normalized;
    \item compute marginal distributions;
    \item determine whether factorization gives independence;
    \item compute conditional distributions when needed;
    \item sketch the support region for continuous problems;
    \item choose appropriate sums or integration bounds;
    \item verify conditional distributions normalize to \(1\);
    \item interpret the final result in terms of dependence.
\end{enumerate}

%%%%%%%%%%%%%%%%%%%%%%%%%%%%%%%%%%%%%%%%%%%%%%%%%%%%%%%%%%%%
\subsection{Common Errors}
\label{subsec:joint-distributions-common-errors}
%%%%%%%%%%%%%%%%%%%%%%%%%%%%%%%%%%%%%%%%%%%%%%%%%%%%%%%%%%%%

\subsubsection{Confusing Joint and Marginal Probability}

The joint quantity

\[
P(X=x,Y=y)
\]

describes simultaneous values.

The marginal quantity

\[
P(X=x)
\]

ignores \(Y\).

%%%%%%%%%%%%%%%%%%%%%%%%%%%%%%%%%%%%%%%%%%%%%%%%%%%%%%%%%%%%
\subsubsection{Forgetting to Sum or Integrate Out a Variable}

Marginal distributions require summing or integrating across all values
of the other variable.

%%%%%%%%%%%%%%%%%%%%%%%%%%%%%%%%%%%%%%%%%%%%%%%%%%%%%%%%%%%%
\subsubsection{Assuming Independence from Marginals Alone}

Knowing \(p_X\) and \(p_Y\) does not determine whether \(X\) and \(Y\)
are independent.

The joint distribution is required.

%%%%%%%%%%%%%%%%%%%%%%%%%%%%%%%%%%%%%%%%%%%%%%%%%%%%%%%%%%%%
\subsubsection{Using Independence Without Checking Factorization}

To claim independence, verify

\[
p_{X,Y}(x,y)=p_X(x)p_Y(y)
\]

or

\[
f_{X,Y}(x,y)=f_X(x)f_Y(y).
\]

%%%%%%%%%%%%%%%%%%%%%%%%%%%%%%%%%%%%%%%%%%%%%%%%%%%%%%%%%%%%
\subsubsection{Ignoring the Shape of the Support}

For continuous joint densities, incorrect support geometry usually
produces incorrect integral limits.

Sketching the region is often essential.

%%%%%%%%%%%%%%%%%%%%%%%%%%%%%%%%%%%%%%%%%%%%%%%%%%%%%%%%%%%%
\subsubsection{Assuming Rectangular Support Implies Independence}

Rectangular support is not enough.

The density must also factor.

%%%%%%%%%%%%%%%%%%%%%%%%%%%%%%%%%%%%%%%%%%%%%%%%%%%%%%%%%%%%
\subsubsection{Confusing Conditional Density with Joint Density}

The conditional density must be normalized with respect to the
conditioning value:

\[
f_{X\mid Y}(x\mid y)
=
\frac{
f_{X,Y}(x,y)
}{
f_Y(y)
}.
\]

%%%%%%%%%%%%%%%%%%%%%%%%%%%%%%%%%%%%%%%%%%%%%%%%%%%%%%%%%%%%
\subsubsection{Conditioning on a Zero-Probability Point Naively}

For continuous variables,

\[
P(Y=y)=0.
\]

Conditional densities are therefore not defined using the elementary
ratio

\[
\frac{P(X\in A,Y=y)}{P(Y=y)}.
\]

They require the density or more advanced conditional-probability
framework.

%%%%%%%%%%%%%%%%%%%%%%%%%%%%%%%%%%%%%%%%%%%%%%%%%%%%%%%%%%%%
\subsubsection{Assuming Zero Correlation Implies Independence}

In general,

\[
\rho_{X,Y}=0
\]

does not imply independence.

The jointly normal case is a special exception.

%%%%%%%%%%%%%%%%%%%%%%%%%%%%%%%%%%%%%%%%%%%%%%%%%%%%%%%%%%%%
\subsubsection{Confusing Marginal and Conditional Independence}

Variables can be independent marginally but dependent conditionally, or
dependent marginally but conditionally independent.

%%%%%%%%%%%%%%%%%%%%%%%%%%%%%%%%%%%%%%%%%%%%%%%%%%%%%%%%%%%%
\subsection{Applications}
\label{subsec:joint-distributions-applications}
%%%%%%%%%%%%%%%%%%%%%%%%%%%%%%%%%%%%%%%%%%%%%%%%%%%%%%%%%%%%

\begin{applicationbox}

\textbf{Statistics.}
Joint distributions model samples, estimators, covariance, regression,
and multivariate data.

\medskip

\textbf{Engineering.}
Multiple uncertain loads, temperatures, stresses, and failure modes are
modeled jointly.

\medskip

\textbf{Finance.}
Asset returns and losses require joint dependence models, not merely
individual marginal distributions.

\medskip

\textbf{Machine Learning.}
Probabilistic graphical models, mixture models, latent-variable models,
and Bayesian networks rely on joint and conditional distributions.

\medskip

\textbf{Biology.}
Joint models describe correlated biological measurements and interacting
population quantities.

\medskip

\textbf{Operations Research.}
Demand, service time, cost, and resource availability may be jointly
random.

\medskip

\textbf{Stochastic Processes.}
A stochastic process is specified through joint distributions of
collections of random variables indexed by time or space.

\medskip

\textbf{Reliability.}
Dependent component lifetimes and shared environmental stresses require
joint distributions.

\end{applicationbox}

%%%%%%%%%%%%%%%%%%%%%%%%%%%%%%%%%%%%%%%%%%%%%%%%%%%%%%%%%%%%
\subsection{Exercises}
\label{subsec:joint-distributions-exercises}
%%%%%%%%%%%%%%%%%%%%%%%%%%%%%%%%%%%%%%%%%%%%%%%%%%%%%%%%%%%%

\begin{exercise}[1]
Define a random vector.
\end{exercise}

\begin{exercise}[1]
Define a joint PMF.
\end{exercise}

\begin{exercise}[1]
State the normalization condition for a joint PMF.
\end{exercise}

\begin{exercise}[1]
Define a marginal PMF.
\end{exercise}

\begin{exercise}[1]
Define a conditional PMF.
\end{exercise}

\begin{exercise}[1]
State the factorization criterion for independence of discrete random
variables.
\end{exercise}

\begin{exercise}[1]
Define a joint PDF.
\end{exercise}

\begin{exercise}[1]
State the normalization condition for a joint PDF.
\end{exercise}

\begin{exercise}[1]
Define the joint CDF.
\end{exercise}

\begin{exercise}[1]
State the factorization criterion for independence using CDFs.
\end{exercise}

\begin{exercise}[2]
Suppose the joint PMF is

\[
\begin{array}{c|cc}
 & Y=0 & Y=1\\
\hline
X=0 & 0.1 & 0.2\\
X=1 & 0.3 & 0.4
\end{array}
\]

Find the marginal PMF of \(X\).
\end{exercise}

\begin{exercise}[2]
For the same table, find the marginal PMF of \(Y\).
\end{exercise}

\begin{exercise}[2]
For the same table, find

\[
P(X=1\mid Y=0).
\]
\end{exercise}

\begin{exercise}[2]
Determine whether \(X\) and \(Y\) are independent in the previous
table.
\end{exercise}

\begin{exercise}[2]
Suppose

\[
f_{X,Y}(x,y)=c
\]

on

\[
0<x<1,
\qquad
0<y<2.
\]

Find \(c\).
\end{exercise}

\begin{exercise}[2]
For the previous joint density, find \(f_X(x)\).
\end{exercise}

\begin{exercise}[2]
For the same density, find \(f_Y(y)\).
\end{exercise}

\begin{exercise}[2]
Determine whether the variables in the previous example are independent.
\end{exercise}

\begin{exercise}[2]
Suppose

\[
f_{X,Y}(x,y)=2
\]

on

\[
0<y<x<1.
\]

Verify that this is a valid density.
\end{exercise}

\begin{exercise}[2]
For the previous density, find the marginal density of \(X\).
\end{exercise}

\begin{exercise}[2]
For the same density, find the marginal density of \(Y\).
\end{exercise}

\begin{exercise}[2]
Find

\[
P(X+Y<1)
\]

for the same joint density.
\end{exercise}

\begin{exercise}[2]
Explain why the triangular support implies dependence.
\end{exercise}

\begin{exercise}[2]
Let \(X\) and \(Y\) be independent fair dice. Find

\[
P(X=Y).
\]
\end{exercise}

\begin{exercise}[2]
For the same dice, find

\[
P(X+Y=8).
\]
\end{exercise}

\begin{exercise}[3]
Prove that the marginal PMF satisfies

\[
p_X(x)=\sum_y p_{X,Y}(x,y).
\]
\end{exercise}

\begin{exercise}[3]
Prove that a conditional PMF sums to \(1\).
\end{exercise}

\begin{exercise}[3]
Prove that if discrete random variables are independent, then

\[
P(X\in A,Y\in B)
=
P(X\in A)P(Y\in B)
\]

for arbitrary subsets \(A,B\).
\end{exercise}

\begin{exercise}[3]
Prove the analogous statement for continuous random variables using
double integrals.
\end{exercise}

\begin{exercise}[3]
Suppose

\[
f_{X,Y}(x,y)
=
c(x+y)
\]

on

\[
0<x<1,
\qquad
0<y<1.
\]

Find \(c\).
\end{exercise}

\begin{exercise}[3]
For the previous density, find both marginals.
\end{exercise}

\begin{exercise}[3]
Determine whether the variables are independent.
\end{exercise}

\begin{exercise}[3]
Find the conditional density

\[
f_{X\mid Y}(x\mid y)
\]

for the previous model.
\end{exercise}

\begin{exercise}[3]
Let

\[
f_{X,Y}(x,y)=8xy
\]

for

\[
0<x<y<1.
\]

Verify normalization and find the marginal densities.
\end{exercise}

\begin{exercise}[3]
For the previous model, calculate

\[
P(X<1/2,\;Y>3/4).
\]
\end{exercise}

\begin{exercise}[3]
Derive the rectangle probability formula from the joint CDF using
inclusion--exclusion.
\end{exercise}

\begin{exercise}[3]
Show that

\[
F_X(x)
=
\lim_{y\to\infty}
F_{X,Y}(x,y).
\]
\end{exercise}

\begin{exercise}[3]
Suppose \(X\) and \(Y\) are independent continuous variables. Derive
the convolution formula for \(X+Y\).
\end{exercise}

\begin{exercise}[3]
Let \(X,Y\) be independent \(\operatorname{Uniform}(0,1)\) random
variables. Find the CDF and density of

\[
M=\max(X,Y).
\]
\end{exercise}

\begin{exercise}[3]
For the same variables, find the CDF and density of

\[
L=\min(X,Y).
\]
\end{exercise}

\begin{exercise}[3]
Give an example of two random variables with the same marginals but two
different joint distributions.
\end{exercise}

\begin{exercise}[3]
Construct a pair of dependent random variables with zero covariance.
\end{exercise}

\begin{exercise}[3]
Explain why i.i.d. random variables are exchangeable.
\end{exercise}

\begin{exercise}[4]
Study copulas and Sklar's theorem.
\end{exercise}

\begin{exercise}[4]
Investigate the Fr\'echet--Hoeffding bounds for joint distribution
functions.
\end{exercise}

\begin{exercise}[4]
Study conditional independence in Bayesian networks.
\end{exercise}

\begin{exercise}[4]
Investigate Simpson's paradox as a consequence of changing conditional
and marginal distributions.
\end{exercise}

\begin{exercise}[4]
Study the multivariate normal distribution and derive its density.
\end{exercise}

\begin{exercise}[4]
Prove that uncorrelated jointly normal random variables are independent.
\end{exercise}

\begin{exercise}[4]
Study Gaussian copulas and compare them with multivariate normal
distributions.
\end{exercise}

\begin{exercise}[4]
Investigate exchangeability and de Finetti's theorem.
\end{exercise}

\begin{exercise}[4]
Study mixture models and latent-variable representations.
\end{exercise}

\begin{exercise}[4]
Investigate regular conditional distributions in measure-theoretic
probability.
\end{exercise}

%%%%%%%%%%%%%%%%%%%%%%%%%%%%%%%%%%%%%%%%%%%%%%%%%%%%%%%%%%%%
\subsection{Conceptual Summary}
\label{subsec:joint-distributions-summary}
%%%%%%%%%%%%%%%%%%%%%%%%%%%%%%%%%%%%%%%%%%%%%%%%%%%%%%%%%%%%

A random vector is

\[
\boxed{
\mathbf{X}
=
(X_1,\ldots,X_n).
}
\]

For discrete variables,

\[
\boxed{
p_{X,Y}(x,y)
=
P(X=x,Y=y).
}
\]

Marginals are obtained by summing:

\[
\boxed{
p_X(x)
=
\sum_y
p_{X,Y}(x,y).
}
\]

Conditional PMFs satisfy

\[
\boxed{
p_{X\mid Y}(x\mid y)
=
\frac{
p_{X,Y}(x,y)
}{
p_Y(y)
}.
}
\]

Discrete independence is equivalent to

\[
\boxed{
p_{X,Y}(x,y)
=
p_X(x)p_Y(y).
}
\]

For continuous variables,

\[
\boxed{
P((X,Y)\in A)
=
\iint_A
f_{X,Y}(x,y)\,dx\,dy.
}
\]

Marginal densities are

\[
\boxed{
f_X(x)
=
\int_{-\infty}^{\infty}
f_{X,Y}(x,y)\,dy,
}
\]

and

\[
\boxed{
f_Y(y)
=
\int_{-\infty}^{\infty}
f_{X,Y}(x,y)\,dx.
}
\]

Conditional densities satisfy

\[
\boxed{
f_{X\mid Y}(x\mid y)
=
\frac{
f_{X,Y}(x,y)
}{
f_Y(y)
}.
}
\]

Continuous independence is equivalent to

\[
\boxed{
f_{X,Y}(x,y)
=
f_X(x)f_Y(y).
}
\]

The joint CDF is

\[
\boxed{
F_{X,Y}(x,y)
=
P(X\leq x,Y\leq y).
}
\]

Independence is also equivalent to

\[
\boxed{
F_{X,Y}(x,y)
=
F_X(x)F_Y(y).
}
\]

Marginals alone do not determine dependence.

Thus,

\[
\boxed{
\text{joint distribution}
=
\text{marginals}
+
\text{dependence structure}.
}
\]

Joint distributions provide the foundation for

\[
\boxed{
\text{covariance},
\quad
\text{correlation},
\quad
\text{multivariate statistics},
\quad
\text{stochastic processes}.
}
\]

%%%%%%%%%%%%%%%%%%%%%%%%%%%%%%%%%%%%%%%%%%%%%%%%%%%%%%%%%%%%
\subsection{Section Connections}
\label{subsec:joint-distributions-connections}
%%%%%%%%%%%%%%%%%%%%%%%%%%%%%%%%%%%%%%%%%%%%%%%%%%%%%%%%%%%%

\chapterconnection{
Joint Distributions extends the one-variable theory of Sections
\ref{sec:random-variables}, \ref{sec:discrete-distributions}, and
\ref{sec:continuous-distributions} to random vectors. The next section,
Expectation and Moments, uses joint distributions to define covariance,
correlation, conditional expectation, and moments of several variables.
Functions of Random Variables develops sums, ratios, maxima, minima,
Jacobians, and convolution in greater depth. Probability Transforms
provides transform-based methods for sums and joint laws. The Laws of
Large Numbers and Central Limit Theorems study large collections of
random variables, while Stochastic Processes can be viewed as families
of random variables whose finite-dimensional behavior is described by
joint distributions.
}

%%%%%%%%%%%%%%%%%%%%%%%%%%%%%%%%%%%%%%%%%%%%%%%%%%%%%%%%%%%%
% SECTION SOURCE TRACKING
%%%%%%%%%%%%%%%%%%%%%%%%%%%%%%%%%%%%%%%%%%%%%%%%%%%%%%%%%%%%

%------------------------------------------------------------
% 7.5 Joint Distributions
%------------------------------------------------------------
%
% Source basis:
%   - Standard multivariate probability theory.
%   - Standard joint distribution theory.
%   - Standard conditional probability and density methods.
%
% Used for:
%   - random vectors;
%   - joint distributions;
%   - joint PMFs;
%   - joint PDFs;
%   - joint supports;
%   - probability tables;
%   - discrete marginal distributions;
%   - continuous marginal densities;
%   - conditional PMFs;
%   - conditional densities;
%   - discrete independence;
%   - continuous independence;
%   - support geometry;
%   - probabilities over planar regions;
%   - joint CDFs;
%   - marginal CDFs;
%   - rectangle probabilities;
%   - joint density recovery;
%   - CDF factorization;
%   - sums and convolution preview;
%   - maxima and minima preview;
%   - covariance preview;
%   - correlation preview;
%   - zero covariance versus independence;
%   - conditional independence;
%   - exchangeability;
%   - i.i.d. random variables;
%   - higher-dimensional distributions;
%   - multivariate normal preview;
%   - mixture distributions;
%   - hierarchical models;
%   - dependence structure;
%   - copula preview;
%   - applications and exercises.
%
% Notes:
%   - Expectation, covariance, and correlation are developed fully in
%     Section 7.6.
%   - Transformations, convolution, and Jacobians are developed more
%     systematically in Section 7.7.
%   - Multivariate normal methods are revisited in Statistics.
%   - Formal conditional probability and regular conditional
%     distributions are developed later in measure-theoretic
%     probability.
%
%%%%%%%%%%%%%%%%%%%%%%%%%%%%%%%%%%%%%%%%%%%%%%%%%%%%%%%%%%%%